\section{Discussion}
The reduction to a permanent takes place after the last rows are chosen. The uniform order gives
the geometric bound even for a fixed solution. Under the adaptive law, the same phenomenon follows
from the survival probabilities of prescribed rows and columns. This distinction explains why the
two bounds on the count can use the same last board even though their histories have different
distributions.

The permanent estimate is uniform over a class with fixed diffusion and gap bounds. Its remainder
tends to zero as the size of the last board grows. The order of limits in Section~\ref{balance}
allows us to use precisely this information. The $O(1/n)$ rate in Conjecture~10.1 of~\cite{pred}
remains open. In addition to the permanent error $\theta_K$, the final comparison involves the bound
$C/K+O_\delta(1/n)$ from cumulative consistency (Corollary~\ref{hist:consistencia}), the cost of exits
by band, of order $e^{-c'K}$, and the exceptional contribution of the last board, of order
$K^5\log(K+1)/n$ (Proposition~\ref{terminal-error}). With these bounds, no choice of $K$ as a
function of $n$ gives the rate: $C/K=O(1/n)$ requires $K\ge cn$, in which case
$K^5\log(K+1)/n$ does not tend to zero. To obtain the rate by this method one would have to improve
these bounds and control the permanent error when choosing the floor $K$, or justify cancellations
among the contributions.

There are two close precedents for a result of this type. Eberhard, Manners and
Mrazović~\cite{emm} obtained the asymptotic equivalent of the number of placements of semiqueens
on the toroidal board by a version of the circle method. Tripathi~\cite{tripathi} identifies the
limiting constant of the partition function of a class of Mallows models with a Fredholm
determinant, and also uses McCullagh's estimate~\cite{mccullagh}.
